\section{The paper: setup and main results}

\begin{frame}{Setting: circuits as tensor networks in the superoperator picture}
    \begin{itemize}
        \item Model of Aharonov, Kitaev and Nisan: a gate is any superoperator $Q:\L(\H^{\otimes a})\ra\L(\H^{\otimes b})$ on a constant number of qubits. Unitaries, channels, measurements and trace-outs are all gates.
        \item Every wire carries an index over $\Pi=\{\ket{b_1}\bra{b_2}\}$, so index dimension is $4$, not $2$. Density operators and gates become tensors:
        \[
            \rho_{\sigma_1\cdots\sigma_a}=\tr\rbcs{\rho\,(\otimes_i\sigma_i)^\dagger},\qquad
            Q_{\sigma_1\cdots\sigma_a,\tau_1\cdots\tau_b}=\tr\rbcs{Q(\otimes_i\sigma_i)\,(\otimes_j\tau_j)^\dagger}.
        \]
        \item \textbf{Circuit graph} $\GC$: one vertex per gate, one vertex per open input/output wire, one edge per wire segment. Bounded gate arity $\Ra$ bounded degree.
        \item Attaching $\ket{x_i}\bra{x_i}$ to inputs and POVM elements $\tau(i)$ to outputs gives $N(\C;x,\tau)$. Contracting it to a rank-0 tensor yields the probability that scenario $\tau$ is observed (Prop.~3.5).
        \item Contraction is order-independent in value, but cost depends strongly on the order: $O(T\exp[O(d)])$ where $d$ is the maximal rank encountered (Prop.~3.6).
    \end{itemize}
\end{frame}

\begin{frame}{Contraction complexity and treewidth}
    \begin{dfn}[Contraction complexity]
        A contraction ordering $\pi$ of $E(G)$ merges endpoints one edge at a time. Its complexity is the maximal degree of any merged vertex. $\cc(G)$ is the minimum over all orderings.
    \end{dfn}
    \begin{prop}[Prop.~4.2]
        $\cc(G)=\tw(G^*)$, where $G^*$ is the line graph of $G$. Given a tree decomposition of $G^*$ of width $d$, a contraction ordering of complexity $\le d$ is obtained in polynomial time.
    \end{prop}
    \begin{lem}[Lemma 4.4, bounded degree]
        For any graph $G$ with maximal degree $\Delta(G)$,
        \[
            \tfrac{1}{2}\rbcs{\tw(G)-1}\;\le\;\tw(G^*)\;\le\;\Delta(G)\rbcs{\tw(G)+1}-1 .
        \]
    \end{lem}
    \begin{itemize}
        \item Edge contraction in $G$ $\Lra$ vertex elimination in $G^*$; the degree of the merged vertex equals the induced width of the eliminated vertex. Hence contraction cost is exactly an elimination-width problem.
        \item Treewidth is NP-hard, but the Robertson--Seymour constant-factor approximation runs in $|V|^{O(1)}\exp[O(\tw)]$, which is within the target budget.
    \end{itemize}
\end{frame}

\begin{frame}{Main theorem}
    \begin{thm}[Thm.~1.1 / 4.6]
        Let $\C$ be a circuit with $T$ gates, input $x$ and measurement scenario $\tau$. The probability that $\tau$ is realized on $\C(\rho_x)$ can be computed deterministically in time
        \[
            T^{O(1)}\exp\sbcs{O(\cc(\GC))}\;=\;T^{O(1)}\exp\sbcs{O(\tw(\GC))}.
        \]
    \end{thm}
    \textbf{Algorithm.}
    \begin{enumerate}
        \item Build $N=N(\C;x,\tau)$.
        \item Tree-decompose the line graph $N^*$ with width $O(\tw(N^*))$ (Robertson--Seymour).
        \item Convert the decomposition into a contraction ordering $\pi$ (Prop.~4.2).
        \item Contract along $\pi$; the final rank-0 tensor is the answer.
    \end{enumerate}
    \begin{rem}
        Gates need not be simulated in temporal order; a gate may even be contracted partially at several stages. Arbitrary gates are allowed, so the result is orthogonal to gate-library restrictions (Gottesman--Knill, matchgates) and to bounded-entanglement results (Vidal, Jozsa--Linden).
    \end{rem}
\end{frame}

\begin{frame}{Corollaries by circuit topology}
    \begin{itemize}
        \item \textbf{Cor.~1.2.} Polynomial-size circuits with $\tw(\GC)=O(\log T)$ are simulable in polynomial time. Quantum formulas (trees, $\tw=1$) recover Roychowdhury--Vatan.
        \item \textbf{Cor.~1.3.} Series-parallel circuits ($\tw\le 2$): polynomial time.
        \item \textbf{Cor.~1.4.} Planar circuits: $\tw=O(\sqrt{T})$, hence $\exp[O(\sqrt T)]$.
        \item \textbf{Cor.~1.5 / Prop.~5.1.} $q$-local-interacting circuits of depth $D$: $T^{O(1)}\exp[O(qD)]$. Proof by an explicit \emph{path decomposition} whose $i$-th bag is the set of gates straddling the qubit cut $i\mid i+1$. Log-depth nearest-neighbour circuits are classically easy.
        \item \textbf{Thm.~1.8.} Any depth-3 circuit with one- and two-qubit gates: polynomial time, deterministically (after removing degree-1 vertices, all degrees are $\le 2$, so $\cc\le 2$).
    \end{itemize}
\end{frame}

\begin{frame}{Barrier at depth 4 and one-way computation}
    \begin{thm}[Thm.~1.7]
        There is a depth-4 circuit on $n$ qubits with one- and two-qubit gates whose treewidth is $\Omega(n)$.
    \end{thm}
    \begin{itemize}
        \item Construction: the Lubotzky--Phillips--Sarnak 3-regular expander $x\mapsto x\pm1,\,x^{-1}$ on $\ZZ_p\cup\{\infty\}$; its edges split into three matchings, each a layer of two-qubit gates. Expanders satisfy $\cc(G)\ge\varepsilon|V|/4$ (Lemma 5.2).
        \item Depth 3 is easy, depth 4 can already be as hard as the worst case: the topological criterion is sharp.
    \end{itemize}
    \begin{thm}[Thm.~1.6 / 6.6]
        A one-way computation on the graph state $\ket{G}$ is simulable by a randomized algorithm in $|V|^{O(1)}\exp[O(\tw(G))]$; deterministically if the computation is \emph{oblivious} (as in Raussendorf--Briegel).
    \end{thm}
    \begin{itemize}
        \item Tools: split each $\Lambda(\sigma_z)$ tensor across the edge so $\cc(G')=O(\cc(G))$; then a degree-$\le3$ \emph{expansion} of $G$ with $\tw(G')\le\tw(G)+1$, realized on the graph state by two oblivious measurements per inserted vertex (Prop.~6.4).
    \end{itemize}
\end{frame}

\begin{frame}{Open questions left by the paper (Section 7)}
    \begin{enumerate}
        \item \textbf{Treewidth reduction by circuit rewriting.} Can the treewidth of a circuit be lowered systematically while preserving the computed function?
        \item \textbf{Treewidth of arithmetic.} Conjecture: modular exponentiation (hence Shor) needs large treewidth. Surprise in follow-ups: the approximate QFT has small treewidth (Aharonov--Landau--Makowsky, Yoran--Short).
        \item \textbf{Composition.} Treewidth can jump under composition (depth 3 vs.\ depth 4). Which circuit classes stay simulable when interleaved with QFT?
        \item \textbf{MBQC resources.} Characterize graph states that are universal vs.\ efficiently simulable. Width parameters are necessary but not sufficient (Bravyi--Raussendorf toric-code states are simulable at high width).
    \end{enumerate}
\end{frame}
